\documentclass{article}
\usepackage[a4paper, margin=1in]{geometry} % Adjust margins
\usepackage{amsmath, amssymb, amsthm}
\usepackage{tikz}
\usepackage{xcolor}
\usepackage{graphicx}
\usepackage{listings}
\usepackage{fancyhdr}
\usetikzlibrary{automata, positioning}
\usepackage{helvet} % Match font to reference
\pagestyle{fancy}
\usepackage{enumitem,xcolor}


% Ensure proper header spacing
\setlength{\headheight}{15pt} % Adjust height of the header
\setlength{\headsep}{10pt} 

% Define Header/Footer
\fancyhf{}
\fancyhead[L]{\textbf{Advanced Topics in Computer Science}} % Left header
\fancyhead[R]{Universal Turing Machine Explainer} % Right header
\cfoot{\thepage} % Page number at center bottom

% Define colors for section and subsection titles
\usepackage{titlesec}
\definecolor{sectioncolor}{rgb}{0.1, 0.3, 0.7} % Dark blue
\titleformat{\section}{\large\bfseries\color{sectioncolor}}{\thesection}{1em}{}
\titleformat{\subsection}{\bfseries\color{sectioncolor}}{\thesubsection}{1em}{}

% Define custom colors for different text types
\usepackage{tcolorbox}
\tcbuselibrary{theorems}
\newcommand{\important}[1]{\textcolor{red}{#1}}
\newtcbtheorem
  []% init options
  {definition}% name
  {Definition}% title
  {%
    colback=blue!5,
    colframe=blue!35!black,
    fonttitle=\bfseries,
  }% options
  {def}% prefix
\newcommand{\note}[1]{\textcolor{green!50!black}{#1}}

% Theorem environment setup
\newtheorem{theorem}{Theorem}
\newtheorem{lemma}{Lemma}
\newtheorem{corollary}{Corollary}

\definecolor{codegray}{gray}{0.9}
\lstdefinestyle{customc}{
  backgroundcolor=\color{codegray},
  language=C,
  basicstyle=\ttfamily\footnotesize,
  keywordstyle=\color{blue},
  commentstyle=\color{green!50!black},
  stringstyle=\color{red}
}

\usepackage{tikz}
\usetikzlibrary{automata, positioning, arrows, matrix, fit, calc, decorations}
\usetikzlibrary{decorations.pathreplacing}

\usepackage{xcolor}
\usepackage{xstring}
\usepackage{geometry}
\usepackage{array}
\usepackage{booktabs}
\usepackage{amsmath}

\definecolor{statecolor}{RGB}{50,50,200}
\definecolor{dststatecolor}{RGB}{20,20,150}
\definecolor{tapesymbolcolor}{RGB}{180,50,50}
\definecolor{writesymbolcolor}{RGB}{220,50,50}
\definecolor{directioncolor}{RGB}{50,150,50}
\definecolor{redtape}{RGB}{220,50,50}
\definecolor{greentape}{RGB}{50,180,50}
\definecolor{purpletape}{RGB}{180,100,200}
\definecolor{highlightcolor}{RGB}{255,240,200}

\geometry{margin=1.5cm}


\usepackage{tikz}
\usetikzlibrary{positioning, fit, decorations.pathreplacing, calc}
\usepackage{xstring}

\newcommand{\dynamicTapeAdvanced}[7][]{%
    % #1: Optional tikzpicture parameters
    % #2: Title for the tape
    % #3: Head position (0-indexed)
    % #4: Cells per line (use 0 for single line)
    % #5: String of 0s and 1s
    % #6: Pointer positions and labels: "pos1:label1,pos2:label2,..."
    % #7: Bracket ranges and labels: "start1-end1:label1,start2-end2:label2,..."
    \begin{tikzpicture}[
        cell/.style={draw, minimum size=0.8cm, anchor=center},
        marked/.style={fill=yellow!30},
        head/.style={draw, red, thick, minimum width=1cm, minimum height=1.2cm},
        baseline={(0,-0.5cm)},  % Consistent baseline for alignment
        pointer/.style={-stealth, thick, blue},
        #1
    ]
        % Title/Header
        \node[align=center, font=\bfseries, anchor=north] at (0,2) {#2};
        
        % Calculate string length
        \StrLen{#5}[\stringlength]
        
        % Determine if wrapping is needed
        \pgfmathtruncatemacro{\cellsperline}{max(#4,\stringlength)}
        \ifnum#4>0
            \pgfmathtruncatemacro{\cellsperline}{#4}
        \fi
        
        % Define vertical spacing between rows (increased from 1.2cm)
        \def\verticalspace{2cm}
        
        % Draw cells from the string
        \foreach \i in {1,...,\stringlength} {
            % Extract the character at position \i
            \StrChar{#5}{\i}[\currentchar]
            
            % Calculate 0-indexed position
            \pgfmathtruncatemacro{\zeroindex}{\i-1}
            
            % Calculate row and column position
            \pgfmathtruncatemacro{\row}{int(\zeroindex/\cellsperline)}
            \pgfmathtruncatemacro{\col}{int(\zeroindex-\row*\cellsperline)}
            
            % Check if this is the head position
            \pgfmathparse{int(\zeroindex==#3)}
            \ifnum\pgfmathresult=1
                %\node[cell, marked] (c\zeroindex) at (\col*0.8cm, -\row*\verticalspace) {\currentchar};
                \node[cell] (c\zeroindex) at (\col*0.8cm, -\row*\verticalspace) {\currentchar};
            \else
                \node[cell] (c\zeroindex) at (\col*0.8cm, -\row*\verticalspace) {\currentchar};
            \fi
        }
    
        
        % Process pointers
        \foreach \pointerspec in {#6} {
            \IfSubStr{\pointerspec}{:}{%
                % Extract position and label
                \StrBefore{\pointerspec}{:}[\pointerpos]
                \StrBehind{\pointerspec}{:}[\pointerlabel]
                
                % Draw the pointer
                \draw[pointer] ($(c\pointerpos.south) + (0,-0.5cm)$) -- (c\pointerpos.south);
                \node[blue, below] at ($(c\pointerpos.south) + (0,-0.5cm)$) {\pointerlabel};
            }{%
                % No label provided, just position
                \pgfmathtruncatemacro{\pointerpos}{\pointerspec}
                
                % Draw the pointer without label
                \draw[pointer] ($(c\pointerpos.south) + (0,-0.5cm)$) -- (c\pointerpos.south);
            }
        }
        
        % Process brackets - with increased vertical space for the bracket label
        \foreach \bracketspec in {#7} {
            \IfSubStr{\bracketspec}{:}{%
                % Extract range and label
                \StrBefore{\bracketspec}{:}[\bracketrange]
                \StrBehind{\bracketspec}{:}[\bracketlabel]
                
                % Extract start and end
                \StrBefore{\bracketrange}{-}[\bracketstart]
                \StrBehind{\bracketrange}{-}[\bracketend]
                
                % Calculate rows for start and end
                \pgfmathtruncatemacro{\startrow}{int(\bracketstart/\cellsperline)}
                \pgfmathtruncatemacro{\endrow}{int(\bracketend/\cellsperline)}
                
                % If on same row, draw a single bracket
                \ifnum\startrow=\endrow
                    \draw[decorate, decoration={brace, amplitude=5pt}, thick] 
                        ([yshift=0.4cm]c\bracketstart.north west) -- 
                        ([yshift=0.4cm]c\bracketend.north east) 
                        node[midway, above=5pt] {\bracketlabel};
                \else
                    % Multiple rows - draw separate brackets for each row
                    % First row: from start to end of line
                    \pgfmathtruncatemacro{\firstrowend}{min((\startrow+1)*\cellsperline-1, \stringlength-1)}
                    \draw[decorate, decoration={brace, amplitude=5pt}, thick] 
                        ([yshift=0.4cm]c\bracketstart.north west) -- 
                        ([yshift=0.4cm]c\firstrowend.north east)
                        node[midway, above=5pt] {\bracketlabel};
                    
                    % Middle rows if any
                    \foreach \r in {\startrow,...,\endrow} {
                        \ifnum\r>\startrow
                            \ifnum\r<\endrow
                                \pgfmathtruncatemacro{\rowstart}{\r*\cellsperline}
                                \pgfmathtruncatemacro{\rowend}{min((\r+1)*\cellsperline-1, \stringlength-1)}
                                \draw[decorate, decoration={brace, amplitude=5pt}, thick] 
                                    ([yshift=0.4cm]c\rowstart.north west) -- 
                                    ([yshift=0.4cm]c\rowend.north east);
                            \fi
                        \fi
                    }
                    
                    % Last row: from start of line to end
                    \pgfmathtruncatemacro{\lastrowstart}{\endrow*\cellsperline}
                    \draw[decorate, decoration={brace, amplitude=5pt}, thick] 
                        ([yshift=0.4cm]c\lastrowstart.north west) -- 
                        ([yshift=0.4cm]c\bracketend.north east);
                \fi
            }{%
                % No label provided
                \StrBefore{\bracketspec}{-}[\bracketstart]
                \StrBehind{\bracketspec}{-}[\bracketend]
                
                % Calculate rows for start and end
                \pgfmathtruncatemacro{\startrow}{int(\bracketstart/\cellsperline)}
                \pgfmathtruncatemacro{\endrow}{int(\bracketend/\cellsperline)}
                
                % If on same row, draw a single bracket
                \ifnum\startrow=\endrow
                    \draw[decorate, decoration={brace, amplitude=5pt}, thick] 
                        ([yshift=0.4cm]c\bracketstart.north west) -- 
                        ([yshift=0.4cm]c\bracketend.north east);
                \else
                    % Multiple rows - draw separate brackets
                    % (Same logic as above but without labels)
                    \pgfmathtruncatemacro{\firstrowend}{min((\startrow+1)*\cellsperline-1, \stringlength-1)}
                    \draw[decorate, decoration={brace, amplitude=5pt}, thick] 
                        ([yshift=0.4cm]c\bracketstart.north west) -- 
                        ([yshift=0.4cm]c\firstrowend.north east);
                    
                    \foreach \r in {\startrow,...,\endrow} {
                        \ifnum\r>\startrow
                            \ifnum\r<\endrow
                                \pgfmathtruncatemacro{\rowstart}{\r*\cellsperline}
                                \pgfmathtruncatemacro{\rowend}{min((\r+1)*\cellsperline-1, \stringlength-1)}
                                \draw[decorate, decoration={brace, amplitude=5pt}, thick] 
                                    ([yshift=0.4cm]c\rowstart.north west) -- 
                                    ([yshift=0.4cm]c\rowend.north east);
                            \fi
                        \fi
                    }
                    
                    \pgfmathtruncatemacro{\lastrowstart}{\endrow*\cellsperline}
                    \draw[decorate, decoration={brace, amplitude=5pt}, thick] 
                        ([yshift=0.4cm]c\lastrowstart.north west) -- 
                        ([yshift=0.4cm]c\bracketend.north east);
                \fi
            }
        }
    \end{tikzpicture}
}
